\chapter{Background and Literature Review} \label{chapter:background}

\section{Material Interfaces} \label{section:material_interfaces}

\label{para:interface_section_intro} % ---------- %

Building on our foundational definitions, we now explore the key functional regions of a material interface, the
engineering properties that govern strength, adhesion, and stability, and the principal physical effects arising at
these junctions, electronic phenomena and energy-transport processes.

We then discuss design parameters for tailoring interfacial behaviour, the role of crystallographic orientation and
polymorph-driven reconstruction, and the generation and prediction of interfacial defects.

Finally, we highlight representative applications, from gate dielectric layers and transition metal dichalcogenides
(TMDC) heterojunctions to energy-storage interfaces, thermoelectric contacts, and optoelectronic devices, and conclude
with current active modelling challenges.

\label{para:interfacial_engineering_properties} % ---------- %

In materials engineering, interfaces govern a wide array of properties including adhesion, corrosion resistance,
mechanical strength, and diffusion kinetics.

The strength and ductility of composite systems or coatings often depend on interfacial bonding and lattice
compatibility.

Tailoring interfacial chemistry or morphology, through methods such as surface functionalisation, alloying, or atomic
layer deposition, is often essential for achieving desired macroscopic performance metrics.

\label{para:electronic_interface_effects} % ---------- %

In electronic devices, interfaces critically influence charge transport, band alignment, and contact resistance.

Semiconductor heterojunctions, metal-semiconductor contacts, and dielectric boundaries are fundamental to the operation
of transistors, diodes, and capacitors.

Localised changes in atomic registry, roughness, or defect density at the interface can significantly alter tunnelling
behaviour, barrier heights, and mobility.

Where interfacial dipoles, charge redistribution, and electrostatic reconstruction result in complex electronic
responses not captured by bulk models.

\label{para:energy_interfaces} % ---------- %

Energy systems are similarly interface-dominated.

In lithium-ion batteries, for example, the solid-electrolyte interphase (SEI) acts as a metastable interfacial layer
that dictates both stability and charge transport.

The performance and longevity of solid-state batteries, fuel cells, and thermoelectric devices can depend on the
engineered properties of interfaces, which affect ionic conductivity, catalytic activity, thermal transport, and
electronic insulation.

In photovoltaics, interfacial band offsets between donor and acceptor layers, or between absorber materials and charge
transport layers, govern charge separation efficiency and open-circuit voltage.

\label{para:interface_design_parameters} % ---------- %

Overall, interfaces are not incidental features but active design parameters.

They serve as sites of emergent phenomena, such as interfacial polarisation, bandgap renormalisation, and phonon
scattering, that cannot be predicted from bulk properties alone.

As such, mastering interfacial science is a gateway to developing high-performance materials and devices.

\label{para:interface_definition} % ---------- %

The structural features that influence interface behaviour include crystallographic orientation, the nature of the phase
boundary, and the presence of defects.

\label{para:orientation_effects} % ---------- %

The relative orientation of crystals on either side of an interface can significantly affect the atomic coordination,
bond continuity, and symmetry properties across the interface.

Epitaxial alignment, where lattice planes across the interface maintain a well-defined relationship, typically results
in low-strain, coherent interfaces.

In contrast, a mismatch in orientation can lead to misfit dislocations or interfacial strain, influencing charge carrier
transport, interfacial states, and mechanical stability.

Orientation also affects the formation energy and electronic reconstruction at the interface, which can lead to
phenomena like interface-induced states or altered band alignments.

\label{para:polymorph_interface_reconstruction} % ---------- %

When the interface forms between distinct structural phases, such as cubic and hexagonal polymorphs, unique bonding
geometries and strain fields arise.

For instance, in layered 2D|3D systems like \ce{HfS2}|\ce{HfO2}, the band alignment changes depending on whether the
materials are stacked or laterally connected, with lateral interfaces showing increased reconstruction due to bonding
mismatch between the different crystalline motifs.

These reconstructions are not merely geometric; they result in different electrostatic and electronic behaviours,
influencing band offsets and charge transfer dynamics.

\label{para:interfacial_defects} % ---------- %

Defects, including vacancies, interstitials, dislocations, and grain boundaries, are often intrinsic to interface
formation and can profoundly modify interfacial properties.

They may serve as scattering centres, recombination sites, or sources of local doping.

In grain boundaries of polycrystalline materials, for example, variations in local stoichiometry and strain can create
regions of localised metallicity or enhanced dielectric response, as shown in systems exhibiting colossal permittivity.

In some systems, such as CCTO (\ce{CaCu3Ti4O12}), interfacial defects and microstructural features have been shown to
dominate the macroscopic dielectric properties.

\label{para:defect_ml_prediction} % ---------- %

In computational studies, predicting the impact of such defects requires careful treatment of their spatial distribution,
charge state, and interaction with surrounding lattice atoms.

Modern approaches like the RAFFLE algorithm allow for machine learning-driven structure prediction that accounts for
void-filling and local bonding environments, enabling more realistic models of defect-laden interfaces.

\label{para:interface_applications} % ---------- %

Depending on their nature, interfaces may enhance or impede material performance.

This subsection surveys several application domains in which understanding and engineering interfacial properties is
essential.

\label{para:gate_dielectric_interface} % ---------- %

In semiconductor transistor devices, the interface between the gate dielectric and the semiconducting channel is a
critical design parameter.

% ---------- %

As traditional \ce{SiO2} dielectrics are replaced by high-$\kappa$ materials such as \ce{HfO2} to meet miniaturisation
demands, understanding charge transfer and band alignment across interfaces like \ce{HfS2}| \ce{HfO2} becomes essential
for ensuring device reliability and performance.

Notably, in such 2D|3D heterostructures, the type and strength of bonding (e.g. van der Waals vs. chemical) directly
influences charge redistribution and can cause significant deviation from band alignments predicted by Anderson's rule,
the vacuum levels of the two materials must be aligned.

\label{para:tmdc_interface_band_alignment} % ---------- %

In layered 2D electronics, especially those using TMDCs, the stacking geometry, layer
orientation, and interfacial disorder all strongly influence the band alignment.

These factors affect carrier dynamics, optical transitions, and switching behaviour.

Traditional models such as Anderson's rule have proven insufficient; corrections using terms such as $\Delta E_\Gamma$ and
$\Delta E_{\text{IF}}$ have been developed to capture hybridisation and interface dipole contributions more accurately.

% todo: add formula for detla E_\Gamma and delta E_IF

These corrections allow for predictive modelling of heterostructure behaviour and support targeted band offset
engineering.

\label{para:energy_storage_interface} % ---------- %

Interfaces also govern performance in energy storage systems.

In all-solid-state batteries, for example, the interface between electrode and solid electrolyte dictates ion mobility,
chemical stability, and space charge layer formation.

Such interfacial processes can lead to resistive interphases or dendrite formation, reducing battery life and safety.

Density Functional Theory (DFT) modelling has shown that interfacial reconstructions and electrostatic gradients can
dominate device behaviour.

\label{para:thermoelectric_interface} % ---------- %

In thermoelectric materials, interfaces are employed to decouple electrical and thermal transport.

Interface-induced phonon scattering suppresses thermal conductivity while preserving electronic conductivity; enhancing
the thermoelectric figure of merit, $ZT$.

Layered TMDCs and heterostructures with nanostructured interfaces exploit this mechanism to improve thermal management.

\label{para:interfaces_active_modelling_challenge} % ---------- %

Across these examples, it is evident that interfaces are not passive boundaries, but active, tunable regions that
critically determine macroscopic functionality.

Their predictive modelling, through DFT, machine learning potentials, or interface-specific structure search tools like
ARTEMIS, remains a central challenge in modern materials design.